\documentclass[12 pt]{article} 
\usepackage{float, amsfonts, amssymb, amsmath, blindtext, multicol, graphicx,
enumitem, listings, hyperref, listings}
\usepackage[font=small,labelfont=bf]{caption}
\graphicspath{ {./images/} }
\lstdefinelanguage{TypeScript}{
  keywords={export, interface, type, string, number, boolean, null, undefined},
  sensitive=true,
  comment=[l]{//},
  morecomment=[s]{/*}{*/},
  morestring=[b]",
  morestring=[b]'
}

\oddsidemargin=-0.5cm
\setlength{\textwidth}{6.5in}
\addtolength{\voffset}{-20pt}
\addtolength{\headsep}{25pt} 



\pagestyle{myheadings}                           	
\markright{Anthony Nguyen, Rowan Lochrin\hfill \today \hfill} 	    


\newcommand{\eqn}[0]{\begin{array}{rcl}}
\newcommand{\eqnend}[0]{\end{array} } 

\newcommand{\qed}[0]{$\square$}

\begin{document}
\noindent \textbf{CSC-696D, Takanori Fujiwara} \par
\noindent \textbf{Final project report}. \\
\noindent \textbf{Novel visualization techniques for git data}. \\
\noindent \textbf{Source Code:
\href{https://github.com/rowan907/CSC-696D-Final-Project}{https://github.com/rowan907/CSC-696D-Final-Project}}.
\par
\noindent \textbf{Hosted project:
\href{https://rowan907.github.io/CSC-696D-Final-Project}{https://rowan907.github.io/CSC-696D-Final-Project}}.
\bigskip

\subsection{Problem Description}
In this project we created several novel visual analytic systems for git
data. The visualizations we propose will help developers quickly gain an understanding of
git repositories. In particular, we seek to give users an understanding of how a
repo changes over time. We will focus on a few key questions.
\begin{enumerate}
    \item Who are the main contributors to a repository?
    \item What were the development priorities for a repository at any given
        time?
    \item Can we identify a time when AI-assisted development became popular on
        a given repository?
    \item What is the best visualization for giving developers an understanding of the
        evolution of project structure over time?
\end{enumerate}
\subsection{Motivation}
Developers often work with large amounts of code and are forced to quickly learn
the customs and conventions of many repositories, this process can be very
time-consuming. By giving developers tools for quickly visualizing the history
of a repository, we can immediately give them an understanding of conventions, key
contributors, and etiquette. \par
Many visualizations have been developed over the years to work with git data,
in this project we developed a handful of new visualizations by first
studying visualizations that others have developed, then either refining them to 
improve their visual clarity or augmenting them with additional data. \par
We developed 3 visualizations for this project. 
\begin{enumerate}
    \item A `stabilized' Sediment Graph based on the work of
        \href{https://github.com/erikbern/git-of-theseus}{Erik Bernhardsson}.
        Bernhardsson proposes a unique way to understand the age of the contents of a
        repo, improved his designs by changing the stacking behavior so as rows
    are added we allow the graph to grow in the positive and negative $y$
        directions. This improves the stability of the graph by keeping
        individual lines of sediment more horizontal. Reducing the visual
        interference caused by the sediment layers traveling diagonally up the
        graph, making there width harder to determine. \par
        We also (based on
        feedback we received from our peers), added to new views to this
        visualization. In addition to dividing the layers of sediment by code
        age we also provided the option to divide the sediment by author,
        supporting both total lines of code by author and number of commits by
        author. These views may be an even better use case then code age, as
        authorship is generally a more important metric.
\item Word cloud style view that allows users to select
        many commits via a time window selection and see a word cloud of
        the commit messages. This will allow users to quickly see the
        development focuses of a repo at a particular time window. This is
        inspired by a view from the paper \textit{Visual Analytics for Understanding Software Development History Through
        Git Metadata Analysis}. We seek to enhance their view by adding a time
        window selection allowing these clouds to be generated for any section
        of the commit history. We also provide users the ability to click a word
        from the word cloud and highlight the commits that include that
        word. Allowing users to quickly see what changes dealt with certain
        topics.
    \item A grouping-based visualization based on a force directed graph that clusters files or directories that
        are often modified together, this will help users understand which parts
        of the codebase are most tightly coupled. We animated this view to show
        the evolution of the repo over time and coupled it with a timeline view
        allowing users to scrub to particular sections of the animation. 
\end{enumerate}
\subsection{Related Research Papers}
\begin{itemize}
\item \href{https://ieeexplore.ieee.org/stamp/stamp.jsp?tp=&arnumber=9222261}{Githru\:
        Visual Analytics for Understanding Software Development History Through
        Git Metadata Analysis} ---
    This work shows an all-in-one visual analytic tool for examining git history
        and provides a lot of great background on git. Many techniques here may
        be useful; in particular, the paper provides a lot of detail on the
        technical aspects of creating visualizations with git metadata. We
        also took inspiration from this paper for our second proposed
        visualization as discussed above.
\item \href{https://mircealungu.com/docs/assets/papers/22-Git-Truck.pdf}{Git-Truck: Hierarchy-Oriented Visualization of
Git Repository Evolution} --- Git-Truck is a hierarchical file-organization-oriented
    visualization tool that represents git repositories as different treemaps.
    Each treemap has its nodes sized by the number of lines of code in the repository, 
    and the color of its nodes can represent a variety of interesting factors
    such as number of commits, top contributors, or even the last change day.
    This is showing the current structure of a repository, rather than the temporal 
    focus we have in this project. Mainly, we emphasize the importance of seeing
    the historical development of features in a repository.
\item \href{https://arxiv.org/pdf/2508.12649}{ChangePrism: Visualizing the
    Essence of Code Changes} --- This paper gives another approach to visualizing
        code changes that relies on semantic analysis. Though this paper's focus
        is more on visualizing individual commits and discrete changes rather
        than the type of broad analysis of larger repositories we did,
        it gives some good background on existing visual analytic techniques.
\end{itemize}
\subsection{Technical Approach}
Our technical approach is at a high level
\begin{enumerate}
    \item Collect \texttt{.git} files containing git meta-data from repos we wish to
        study. 
    \item Use a pre-processing script to execute git commands to query this data
        and create a \texttt{.json} file for each repo containing all commits we
        want to analyze and any useful meta-data like branch, parent, files
        changed etc.
    \item Load our \texttt{.json} files at run time when the users selects a
        particular repo they want to analyze.
    \item Build our visualizations from the processed data.
\end{enumerate}
Our pre-processing script lives inside the project
(\texttt{app/scripts/preprocessgit.mjs}) and can be triggered manually
with \texttt{npm run preprocess:git}. The script uses JavaScript's \texttt{exec}
function to run git with the flag \texttt{--git-dir=[.git file]}, this command
allows git to behave for most commands as if it's inside the project we're
targeting despite only having access to the metadata. This saves us having to
pull the whole project we wish to study. \par
We then run a \texttt{git log} to get all commit hashes across all branches (for
performance reasons we cap this at the 4500 most recent commits when more are
available). Finally, we run git commands on every individual commit to get
information like predecessor and branch. The information we pull from each
commit is bellow and is used to support all three visualizations. \par

\noindent\begin{minipage}{\linewidth}
\begin{multicols}{2}
\begin{lstlisting}[language=TypeScript]
type Commit {
  hash: string;
  parents: Commit[];
  children: Commit[];
  author: string;
  date: string;
  subject: string;
  branch: string | null;
  merge?: string | null;
  files?: CommitFile[];
}
\end{lstlisting}
\columnbreak
\begin{lstlisting}[language=TypeScript]
type CommitFile {
  additions: number;
  deletions: number;
  path: string;
}
\end{lstlisting}
\end{multicols}
\end{minipage}
\\
\\

We then include all json files in the \texttt{public} directory of our frontend
app we configure React to serve all files in this directory, so that at runtime
we can query them as needed. We do this because these files are fairly large and not
bundling all of them with our app directly reduces load time. \footnote{If you're
curious you can see the raw json files by navigating to
https://rowan907.github.io/CSC-696D-Final-Project/flask_commits.json} \par

 % Need to add more details about how the actual visualizations are implemented.

\par To speed up implementation we used a few AI tools primary Claude
Haiku/Sonnet/Opus 4.x, as well as Qwen 3.6 running locally for smaller changes and
refactors. \par
\subsection{Technical Challenges}
We initially tried to get all
the information we'd need from the git log command however this involved manually parsing
commit messages to try and infer things like merges which was not completely
accurate. In the end we settled on the approach of getting the commit hashes
from the log and running a git command for every commit to get the additional
data we needed. This
approach was much slower but more reliable, and since the processing script only
needs to run once the slowdown was acceptable.



\subsection{Selected Dataset(s)}
We selected 4 medium size repos to examine in this project.
\begin{itemize}
    \item Flask
    \item Click
    \item Httpx
    \item Requests
\end{itemize}
\subsection{Expected results and evaluation methodology}
In this project, we seek to create a web application to demonstrate and explain
our visual analytic systems. It will feature the visualizations we've created, 
a short text description of the visualizations—explaining works inspiring the
visualization, the types of analytical questions the visualization seeks to
answer and the strengths and weaknesses of our approach. The app will also feature
a selector that will allow you to view any git repository with any of the
aforementioned visualization techniques. We will then gather feedback from
professional software engineers on the usefulness of our visualizations and
applicability of our tools to their current development workflow.

% evaluation examples include: case studies, user studies, quantitative evaluation



\end{document}
